\section{결론}
\label{sec:conclusion}

기본 CoC는 주행 장면에서 선택한 행동의 이유와 목표를 설명하지만, 같은 목표를 달성하는 여러 궤적 중 어느 범위까지 허용되는지는 충분히 나타내지 못할 수 있다. 본 논문은 이러한 허용 조건을 보완하기 위해 기본 CoC와 실행 가드를 하나의 문맥으로 결합한 \sccoc{}를 제안하였다. 또한 같은 장면, 행동 목표 및 궤적 후보를 유지한 채 실행 가드만 바꾸는 쌍대 계약 평가를 구성하여, 모델이 장면 유형이나 행동 목표가 아니라 실제 가드의 차이를 읽고 궤적을 선택하는지 확인하였다.

Qwen3-VL-2B-Instruct를 미세 조정한 시간 조건 실험에서 요구사항 전용형 CoC의 가드 위반률은 $0.319\pm0.052$였으나, CoC 통합형에서는 $0.000\pm0.000$으로 감소하였다. 이때 안전한 목표 완료율과 쌍대 계약 정확도는 모두 $1.000$을 유지하여, 위반 감소가 단순히 항상 정지하는 보수적 선택에서 나온 결과가 아님을 확인하였다. 이러한 결과는 정지, 출발 및 차선 변경을 포함한 표준 평가에서도 일관되었다. 다만 어려운 표현 평가에서는 위반률이 $0.010\pm0.010$, 쌍대 계약 정확도가 $0.896\pm0.127$이었고, 오류는 주로 m와 cm처럼 단위 표현이 바뀐 조건에 집중되었다. 따라서 작은 VLM이 행동 목표를 유지하면서 실행 가드를 학습할 수는 있지만, 새로운 수치 표현까지 항상 정확하게 해석한다고 볼 수는 없다.

본 연구의 결과는 합성 장면과 미리 정한 궤적 후보를 사용한 작은 VLM의 선택 실험에 한정되며, Alpamayo-R1의 연속 궤적이나 실제 차량의 안전성을 직접 입증하지 않는다. 보조 폐루프 실험에서도 동적 계약과 차단장치를 함께 사용하면 위반이 줄었지만, 계약 갱신이 늦어지면 다시 증가하여 학습된 가드만으로 안전을 보장할 수 없음을 확인하였다. 실제 적용을 위해서는 자연어 가드를 단위가 표준화된 구조화 조건으로 변환하고, 최신 센서 상태를 이용하는 독립적인 궤적 검증기 및 차단장치와 결합해야 한다. 향후에는 이 구조를 대형 자율주행 VLM의 연속 궤적에 적용하고, 실제 센서 증거와 충돌 위험 지표를 이용해 효과를 검증할 필요가 있다.
